\section{Cells and parity}\label{section: cells}

Fix $q$ and the modules $\M^{+} \defeq \M^{+}_q$, $\M^{-} \defeq \M^{-}_q$ of \Cref{proposition: local pair}.
For a bit $\theta \in \set{0, 1}$, let $\M^{(1)} \defeq \M^{+}$ and $\M^{(0)} \defeq \M^{-}$.
So the bit of a module is the truth value of $\Kt(u, v)$.

\begin{definition}[Cell models]\label{definition: cells}
  Let $w = w_1 \cdots w_n \in \set{0,1}^{+}$.
  The valuation $\C(w)$ consists of the following vertices, ordered from left to right:
  \[
    p^{0}_{0} < p^{0}_{1} < \underbrace{u_1 < \cdots < v_1}_{\text{copy of } \M^{(w_1)}} < p^{1}_{0} < p^{1}_{1} < \underbrace{u_2 < \cdots < v_2}_{\text{copy of } \M^{(w_2)}} < \cdots < v_n < p^{n}_{0} < p^{n}_{1}.
  \]
  The relations $\vr, \vc, \vb$ are the union of the relations of the copies of the modules.
  The vertices $p^{i}_{0}, p^{i}_{1}$ are called \emph{ports}.
  Cell $i$ consists of the module $\range{u_i, v_i}$ together with its input ports $p^{i-1}_{0}, p^{i-1}_{1}$ and output ports $p^{i}_{0}, p^{i}_{1}$.
\end{definition}

No input pair leaves a module, and every input pair is strict.
So $\C(w)$ is an $\RSUB$ valuation with irreflexive $\vr, \vc, \vb$.

\begin{lemma}\label{lemma: cells pathwidth}
  $\pw(\C(w)) = 2$.
\end{lemma}
\begin{proof}
  We concatenate path decompositions.
  For each $i$, take the bags $\set{p^{i-1}_0, p^{i-1}_1}$ and $\set{p^{i-1}_1, u_i}$.
  Then take the bags of \Cref{lemma: module shape} for the $i$-th module, and finally the bag $\set{v_i, p^{i}_0}$.
  At the end we add the bag $\set{p^{n}_0, p^{n}_1}$.
  Every vertex occurs in an interval of bags, and every edge is covered, since the ports are joined only by successor pairs.
  All bags have at most three elements.
  The path from $u_1$ to $v_1$ together with the edge $\set{u_1, v_1}$ forms a cycle, so the pathwidth is at least two.
\end{proof}

\begin{figure}[t]
  \centering
  \begin{tikzpicture}[x=1cm,y=1cm,
    port/.style={circle,draw,fill=white,inner sep=0pt,minimum size=5pt},
    endpt/.style={circle,draw,thick,inner sep=0.5pt,minimum size=11pt,font=\scriptsize},
    harr/.style={-{Stealth[length=4pt]},thick},
    succ/.style={gray!60,thin}]
  % ---------- (a) order of the vertices in one cell ----------
  \node[font=\small] at (-0.7,0.35) {(a)};
  \draw[succ] (0,0) -- (7.0,0);
  \node[port] (a0) at (0,0) {}; \node[port] (a1) at (0.7,0) {};
  \node[endpt] (u) at (1.4,0) {$u$};
  \fill[gray!15] (1.75,-0.18) rectangle (4.95,0.18);
  \node[font=\scriptsize] at (3.35,0) {interior of the module};
  \node[endpt] (v) at (5.6,0) {$v$};
  \node[port] (b0) at (6.3,0) {}; \node[port] (b1) at (7.0,0) {};
  \node[font=\scriptsize,below=1pt] at (a0.south) {$p^{i-1}_0$};
  \node[font=\scriptsize,below=1pt,xshift=-2pt] at (a1.south) {$p^{i-1}_1$};
  \node[font=\scriptsize,below=1pt] at (b0.south) {$p^{i}_0$};
  \node[font=\scriptsize,below=1pt] at (b1.south) {$p^{i}_1$};
  \draw[red!70!black,-{Stealth[length=4pt]}] (u) to[bend left=22] node[above,font=\scriptsize]{$\vb$} (v);
  \draw[decorate,decoration={brace,amplitude=3pt}] (0,0.25) -- (1.25,0.25) node[midway,above=2pt,font=\scriptsize]{$\St^{2}$};
  \draw[decorate,decoration={brace,mirror,amplitude=3pt}] (0.7,-0.62) -- (1.4,-0.62) node[midway,below=2pt,font=\scriptsize]{$\St$};
  \draw[decorate,decoration={brace,amplitude=3pt}] (5.75,0.25) -- (7.0,0.25) node[midway,above=2pt,font=\scriptsize]{$\St^{2}$};
  \draw[decorate,decoration={brace,mirror,amplitude=3pt}] (5.6,-0.62) -- (6.3,-0.62) node[midway,below=2pt,font=\scriptsize]{$\St$};
  % ---------- (b) H-edges of one cell, two tracks ----------
  \begin{scope}[shift={(0,-2.9)}]
    \node[font=\small] at (-0.7,1.6) {(b)};
    \foreach \x/\bit/\K in {0/1/\Kt, 4.6/0/\Nt}{
      \begin{scope}[shift={(\x,0)}]
        \node[port] (l0) at (0,0) {}; \node[port] (l1) at (0,1.2) {};
        \node[port] (r0) at (2.6,0) {}; \node[port] (r1) at (2.6,1.2) {};
        \node[font=\scriptsize,left=1pt] at (l0.west) {$p^{i-1}_0$};
        \node[font=\scriptsize,left=1pt] at (l1.west) {$p^{i-1}_1$};
        \node[font=\scriptsize,right=1pt] at (r0.east) {$p^{i}_0$};
        \node[font=\scriptsize,right=1pt] at (r1.east) {$p^{i}_1$};
        \ifnum\bit=1
          \draw[harr] (l0) -- node[pos=0.22,below right=-1pt,font=\scriptsize]{$\St^{2}\compo\Kt\compo\St^{2}$} (r1);
          \draw[harr] (l1) -- node[pos=0.22,above right=-1pt,font=\scriptsize]{$\St\compo\Kt\compo\St$} (r0);
          \node[font=\scriptsize] at (1.3,-0.55) {$w_i = 1$: $\Kt(u_i, v_i)$, swap};
        \else
          \draw[harr] (l0) -- node[below,font=\scriptsize]{$\St^{2}\compo\Nt\compo\St$} (r0);
          \draw[harr] (l1) -- node[above,font=\scriptsize]{$\St\compo\Nt\compo\St^{2}$} (r1);
          \node[font=\scriptsize] at (1.3,-0.55) {$w_i = 0$: $\Nt(u_i, v_i)$, identity};
        \fi
      \end{scope}
    }
  \end{scope}
  % ---------- (c) two-track view of C(101) ----------
  \begin{scope}[shift={(0,-5.6)}]
    \node[font=\small] at (-0.7,1.2) {(c)};
    \foreach \i in {0,1,2,3}{
      \node[port] (q\i0) at (2.2*\i,0) {};
      \node[port] (q\i1) at (2.2*\i,0.9) {};
      \node[font=\scriptsize,below=1pt] at (q\i0.south) {$p^{\i}_0$};
      \node[font=\scriptsize,above=1pt] at (q\i1.north) {$p^{\i}_1$};
    }
    \foreach \i/\j/\b in {0/1/1,1/2/0,2/3/1}{
      \ifnum\b=1
        \draw[harr] (q\i0) -- (q\j1); \draw[harr] (q\i1) -- (q\j0);
      \else
        \draw[harr] (q\i0) -- (q\j0); \draw[harr] (q\i1) -- (q\j1);
      \fi
      \node[font=\scriptsize] at (2.2*\i+1.1,-0.55) {$w_{\j}=\b$};
    }
    \draw[very thick,orange!80!black,opacity=0.5] (q00.center) -- (q11.center) -- (q21.center) -- (q30.center);
    \node[font=\scriptsize,anchor=west,align=left] at (7.0,0.45) {$\Psit(p^0_0,p^3_0)$ holds,\\ $\Psit(p^0_0,p^3_1)$ fails};
  \end{scope}
\end{tikzpicture}

  \caption{(a) The order of the vertices in cell $i$; the braces show the shifts by $\St$ and $\St^{2}$ used in $\Ht$.
    (b) The edges of $\Ht$ in cell $i$, drawn on two tracks.
    A module in $\Kt$ (bit $1$) swaps the ports, and a module in $\Nt$ (bit $0$) keeps them.
    (c) The edges of $\Ht$ in $\C(101)$.
    The highlighted $\Ht$-path from $p^{0}_{0}$ uses two swaps and ends in $p^{3}_{0}$.
    In general, $\Psit(p^{0}_{0}, p^{n}_{1})$ holds iff the number of swaps is odd.}
  \label{figure: cells}
\end{figure}

\begin{lemma}[Parity]\label{lemma: parity}
  In $\C(w)$, the relation $\Ht$ is exactly
  \[
    \set{(p^{i-1}_{j}, p^{i}_{j \oplus w_i}) \mid i \in \range{n},\ j \in \set{0, 1}},
  \]
  where $\oplus$ is addition modulo two.
  Consequently, $\Psit(p^{0}_{0}, p^{n}_{1})$ holds in $\C(w)$ iff the number of $1$s in $w$ is odd.
\end{lemma}
\begin{proof}
  The relation $\vb$ consists of the pairs $(u_i, v_i)$.
  By \Cref{lemma: locality}, $\Kt(u_i, v_i)$ is evaluated in the $i$-th module, so it holds iff $w_i = 1$.
  Hence $\Kt = \set{(u_i, v_i) \mid w_i = 1}$ and $\Nt = \set{(u_i, v_i) \mid w_i = 0}$.
  The only vertex $x$ with $\St^{2}(x, u_i)$ is $p^{i-1}_0$, and the only vertex with $\St(x, u_i)$ is $p^{i-1}_1$.
  Similarly, the only $y$ with $\St(v_i, y)$ is $p^{i}_0$, and the only $y$ with $\St^{2}(v_i, y)$ is $p^{i}_1$.
  Reading off \eqref{equation: Psi} gives the description of $\Ht$; see \Cref{figure: cells}(a,b).
  Every $\Ht$-path from $p^{0}_{0}$ passes through one port of each layer $i$.
  It changes the index exactly at the cells with $w_i = 1$.
\end{proof}
